\documentclass[11pt]{amsart}
\usepackage{geometry}                % See geometry.pdf to learn the layout options. There are lots.
\geometry{letterpaper}                   % ... or a4paper or a5paper or ... 
%\geometry{landscape}                % Activate for for rotated page geometry
%\usepackage[parfill]{parskip}    % Activate to begin paragraphs with an empty line rather than an indent
\usepackage{graphicx}
\usepackage{amssymb}
\usepackage{epstopdf}
\DeclareGraphicsRule{.tif}{png}{.png}{`convert #1 `dirname #1`/`basename #1 .tif`.png}

\title{PS 4.28}
%\date{}                                           % Activate to display a given date or no date

\begin{document}
\maketitle
%\section{}
%\subsection{}
 \noindent Dada una matriz cuadrada MAT (M x N), construya un diagrama de flujo que determine si la misma puede considerarse como un cuadrado mágico. Un cuadrado mágico es aquel en que la suma de las filas, columnas y diagonales siempre tiene el mismo valor. \newline

 \noindent Dato: $A[1..N, 1..N]$ (arreglo bidimensional de tipo entero, $1 \leq N \leq 50$). 



\end{document}  